The practical possibility of recovering a temporally matching 2019 specimen from extant colonies has already reached zero.

This is not a rhetorical flourish or an administrative cutoff. It is a direct, irreversible consequence of molecular evolution operating under ordinary biological conditions. The argument rests on the fundamental properties of viral population genetics, the continuous action of molecular clocks, and the simple arithmetic of elapsed calendar time. Once those elements are stated clearly, the conclusion follows with the force of a physical constraint rather than a matter of incomplete sampling effort.

The continuous operation of molecular clocks
RNA viruses such as sarbecoviruses accumulate nucleotide substitutions, insertions, deletions, and recombination events at measurable rates while they replicate in living hosts. The molecular clock for coronaviruses is typically on the order of \(10^{-3}\) to \(10^{-4}\) substitutions per site per year in the structural genes, with higher rates possible in non-structural regions under certain selective regimes. These rates are not theoretical constructs invented for pandemic debates; they are the same clocks used to date the divergence of RaTG13, the BANAL viruses, RmYN02, and every other published SARS-related sequence. They are calibrated against known sampling dates and are routinely applied in Bayesian phylogenetic frameworks (BEAST and related software) across the literature.

Because the clock never stops while a lineage is circulating, any viral population that existed in 2019 and continued to replicate in bats or intermediate hosts through 2020, 2021, 2022, 2023, 2024, 2025, and into 2026 has necessarily experienced six to seven additional years of mutational accumulation. A specimen collected from those populations in August 2026 therefore belongs to a descendant generation. It cannot be genetically identical—or even sufficiently proximal—to the precise genome that would have been required as the immediate pre-spillover ancestor in late 2019. Evolution does not pause, reverse, or hold a lineage in stasis simply because investigators later decide they need a pristine temporal match.

Temporal irreversibility applied to the proximal-ancestor problem
The defining molecular feature under discussion is the precise polybasic furin cleavage site (PRRA insertion with the CGG-CGG codon pair) that distinguishes SARS-CoV-2 from all previously characterized sarbecoviruses recovered from the intensively surveyed southern Chinese and Southeast Asian reservoirs. No natural specimen carrying that exact configuration has been documented from the field work that ran through May 2019 or from the subsequent years of sampling. The nearest relatives—RaTG13 (≈96.2 % identity), the BANAL series from Laos (≈96.8 % overall), and other regional sarbecoviruses—lack the FCS motif entirely.

Suppose, for the sake of argument, that an undetected natural lineage carrying the exact 2019 configuration had been circulating in some unsampled colony or trade chain in 2019. That lineage would have continued to evolve. By 2026 the descendants of that lineage would already differ from their 2019 progenitors by dozens to hundreds of additional substitutions and possible recombination events, depending on the exact substitution rate and generation time in the reservoir. Recovering a modern specimen and claiming it represents the missing 2019 proximal ancestor would require either (a) that the lineage had somehow ceased evolving for six-plus years—an impossibility under known biology—or (b) that investigators could somehow “rewind” the molecular clock. Neither option exists. Living reservoirs supply only contemporary material; they cannot furnish a frozen 2019 genome.

This is why the practical possibility has already reached zero. The window in which a temporally matching specimen could have been recovered closed with the passage of calendar time itself. Future expeditions into the same karst systems, the same Rhinolophus colonies, or the same wildlife-trade networks can recover closer overall relatives, additional recombinants, or even independent furin-like motifs in distant lineages (as illustrated by certain South American bat coronaviruses). What they cannot recover is a virus that stopped evolving in 2019 and waited patiently for discovery.

Why intensified future sampling cannot reopen the window
The intensity of the pre-2019 sampling network is relevant precisely because it demonstrates that the research apparatus was capable of recovering diverse SARS-related sequences when they were present. The Year-5 RPPR for NIH Grant R01AI110964 documents the scale of that effort: more than 16,000 individual bats, serological surveys across multiple provinces, phylogenetic reconstruction of full-length N genes, and experimental characterization of ACE2 usage. Subsequent independent surveys have continued to expand the known diversity. Yet the critical combination—a natural proximal ancestor already possessing the exact PRRA + CGG-CGG configuration—has remained absent.

Each additional year of negative results further reduces the probability that the exact form was simply overlooked in the surveyed reservoirs. Simultaneously, the evolutionary distance between any living population and the 2019 emergence window continues to increase. The two trends move in the same direction: declining probability of prior oversight and rising impossibility of temporal recovery. There is no compensatory mechanism—neither deeper sequencing of existing colonies nor expansion into previously unsampled caves—that can restore a pristine 2019 genome once evolutionary time has continued.

Structural implications for natural-proximal hypotheses
Pure natural-proximal accounts must therefore abandon any hope of recovering a temporally authentic specimen from extant colonies. They are left with two residual options, both of which multiply unobserved entities:

Entirely unreported or undersampled local populations and trade chains that somehow escaped the documented surveillance network yet managed to generate and maintain the exact FCS-bearing lineage long enough for it to reach Wuhan, after which the lineage either went extinct or continued evolving beyond recognition; or  
An immediate, low-probability multi-species cascade (direct bat-to-human or rapid serial jumps through intermediate hosts) that left no recoverable intermediate specimen and produced a virus already optimally adapted for human transmission upon first detection.

Both options are formally possible. Neither is supported by a verified natural specimen. Both stand in contrast to the research-related pathway, which does not require a hidden natural proximal ancestor circulating outside the documented field-to-lab continuum. Samples moved from high-risk bat sites into laboratory systems already engaged in culture, sequencing, ACE2 assays, and spike-protein experimentation; the geographic and temporal terminus of that continuum coincides with the single-point origin signature first detectable in Wuhan.

The evolutionary-clock barrier is therefore not neutral. It permanently tightens the requirements on natural-proximal narratives while leaving the alternative pathway unconstrained by the same missing-specimen and temporal-irreversibility problems. The evidence does not prove a laboratory-associated introduction. It simply records that the practical possibility of recovering a temporally matching 2019 natural specimen from living colonies has already reached zero, and that this zero is a physical consequence of molecular evolution rather than a temporary gap in fieldwork.
